\documentclass[a4paper,10pt,brazil]{article}
\usepackage{provastyle}
\usepackage{menukeys}

\begin{document}
\makeexamtitlepage{UFOP}{BCC222: Programação Funcional}{José Romildo}{2026/1}{Primeira Prova}{10}
\begin{multicols}{2}
\questionTitle{1}{Análise de assinatura de tipo}
\begin{flushright}
  \footnotesize
  \textit{\textbf{4: Tipos de Dados}}\\
  \textit{c04-assinatura-analise-01}\\

\end{flushright}
Qual das seguintes definições de função corresponde à assinatura de tipo \haskellinline{String -> Int -> Bool}?

\begin{enumerate}
  \item \begin{haskellcode}
verifica s = s == "OK"
\end{haskellcode}

  \item \begin{haskellcode}
verifica s n = length s > n
\end{haskellcode}

  \item \begin{haskellcode}
verifica s n = s ++ show n
\end{haskellcode}

  \item \begin{haskellcode}
verifica n s = n > 0 && s == "OK"
\end{haskellcode}

  \item \begin{haskellcode}
verifica b s = not b && null s
\end{haskellcode}


\end{enumerate}

\questionTitle{2}{Enum e Bounded não são equivalentes}
\begin{flushright}
  \footnotesize
  \textit{\textbf{7: Polimorfismo}}\\
  \textit{cap07-hierarquia-enum-bounded}\\

\end{flushright}
Qual afirmação distingue corretamente as classes \haskellinline{Enum} e \haskellinline{Bounded}?
\begin{enumerate}
  \item \haskellinline{Bounded} permite usar \haskellinline{succ}, enquanto \haskellinline{Enum} fornece \haskellinline{minBound}.
  \item \haskellinline{Enum} descreve tipos que podem ser enumerados com operações como \haskellinline{succ} e \haskellinline{pred}; \haskellinline{Bounded} descreve tipos com \haskellinline{minBound} e \haskellinline{maxBound}.
  \item Todo tipo instância de \haskellinline{Enum} é necessariamente instância de \haskellinline{Bounded}.
  \item As duas classes pertencem apenas à hierarquia de \haskellinline{Fractional}.
  \item \haskellinline{Enum} e \haskellinline{Bounded} são a mesma classe com nomes diferentes.

\end{enumerate}

\questionTitle{3}{A sintaxe de múltiplas linhas}
\begin{flushright}
  \footnotesize
  \textit{\textbf{2: Ambiente de desenvolvimento}}\\
  \textit{c02-fluxo-trabalho-03}\\

\end{flushright}
Ao tentar copiar uma definição funcional de 4 linhas de um PDF e colar diretamente no \textbf{prompt} do GHCi, o terminal costuma gerar erros acusando que identificadores não estão em escopo, pois avalia cada linha independentemente. 

Qual é a solução mais segura, utilizando delimitadores, para que o GHCi aceite todo o bloco colado como uma única construção sintática?

\begin{enumerate}
  \item Inserir o comando de avaliação em lote \texttt{:batch} na linha imediatamente anterior.

  \item Envolver a expressão utilizando parênteses duplos \texttt{((} no início e \texttt{))} no fim.

  \item Utilizar aspas triplas \texttt{"""} tanto no topo quanto na base do trecho copiado.

  \item Involucrar o bloco copiado digitando \texttt{:\{} antes de colar, e \texttt{:\}} ao finalizar.

  \item Adicionar barras duplas \haskellinline{//} no fim de cada linha para avisar o interpretador de que o comando continua na próxima.


\end{enumerate}

\questionTitle{4}{Regra de nomenclatura de tipos}
\begin{flushright}
  \footnotesize
  \textit{\textbf{4: Tipos de Dados}}\\
  \textit{c04-conceito-nomes-01}\\

\end{flushright}
Qual é a regra sintática obrigatória para nomes de tipos em Haskell?

\begin{enumerate}
  \item Devem ser escritos inteiramente em maiúsculas.

  \item Não podem conter caracteres especiais.

  \item Devem começar com uma letra minúscula.

  \item Devem conter pelo menos um dígito.

  \item Devem começar com uma letra maiúscula.


\end{enumerate}

\questionTitle{5}{Avaliação lazy e argumento não usado}
\begin{flushright}
  \footnotesize
  \textit{\textbf{3: Expressões e definições}}\\
  \textit{c03-layout-avaliacao-03}\\

\end{flushright}
Considere a expressão:

\begin{haskellcode}
(\x -> 1) (div 1 0)
\end{haskellcode}

O que acontece em Haskell, considerando a avaliação \emph{lazy}?

\begin{enumerate}
  \item A expressão é sintaticamente inválida, pois uma lambda não pode ser aplicada diretamente a um argumento.

  \item A expressão avalia para \haskellinline|1|, pois o argumento \haskellinline|div 1 0| não é necessário para produzir o resultado da função.

  \item A expressão avalia para a própria função \haskellinline|\x -> 1|, sem aplicar o argumento.

  \item A expressão sempre produz erro de divisão por zero antes de a função ser aplicada, como ocorreria em avaliação estrita.

  \item A expressão avalia para \haskellinline|0|, pois Haskell define \haskellinline|div 1 0| como zero em contextos não usados.


\end{enumerate}

\questionTitle{6}{Soma produzida por leitura até sentinela}
\begin{flushright}
  \footnotesize
  \textit{\textbf{10: Ações de E/S recursivas}}\\
  \textit{c10-rastreamento-01}\\

\end{flushright}
A função abaixo é executada, e o usuário digita \haskellinline|3, 3, 3, 3, 0|, um valor por linha.

\begin{haskellcode}
lerESomar :: IO Integer
lerESomar = do
  n <- readLn
  if n == 0
    then return 0
    else do
      somaResto <- lerESomar
      return (n + somaResto)
\end{haskellcode}

Qual valor é produzido por \haskellinline|lerESomar|?

\begin{enumerate}
  \item 12

  \item 13

  \item 3

  \item 4

  \item 0


\end{enumerate}

\questionTitle{7}{A primeira linguagem funcional prática}
\begin{flushright}
  \footnotesize
  \textit{\textbf{1: Paradigmas de programação}}\\
  \textit{c01-historia-evolucao-02}\\

\end{flushright}
Criada na virada da década de 1950 por John McCarthy no MIT, qual linguagem de programação é historicamente consagrada como o primeiro esforço prático para incorporar o paradigma funcional, introduzindo conceitos inovadores à época, como a recursão nativa e o \emph{Garbage Collection}?

\begin{enumerate}
  \item ML (MetaLanguage)

  \item COBOL

  \item Miranda

  \item Lisp

  \item Fortran


\end{enumerate}

\questionTitle{8}{A distinção entre GHC e GHCi}
\begin{flushright}
  \footnotesize
  \textit{\textbf{2: Ambiente de desenvolvimento}}\\
  \textit{c02-ecossistema-ferramentas-03}\\

\end{flushright}
O conjunto de ferramentas do Glasgow Haskell Compiler fornece dois componentes essenciais: o \texttt{ghc} e o \texttt{ghci}. Qual é a distinção primária entre eles durante o desenvolvimento?

\begin{enumerate}
  \item O \texttt{ghc} é o compilador de linha de comando que gera binários executáveis, enquanto o \texttt{ghci} é o ambiente interativo (\emph{REPL}) para avaliação em tempo real.

  \item O \texttt{ghc} é utilizado exclusivamente no sistema operacional Linux, enquanto o \texttt{ghci} é a versão portada para o Windows.

  \item Não há diferença funcional; ambos são atalhos distintos de terminal que iniciam exatamente a mesma ferramenta de \emph{build}.

  \item O \texttt{ghc} gerencia a formatação do código no editor de texto, ao passo que o \texttt{ghci} executa os testes de unidade automatizados.

  \item O \texttt{ghc} traduz funções puras matemáticas, enquanto o \texttt{ghci} é reservado apenas para compilar funções que realizam operações de entrada e saída (I/O).


\end{enumerate}

\questionTitle{9}{Modelo de E/S em uma linguagem pura}
\begin{flushright}
  \footnotesize
  \textit{\textbf{8: Programas interativos}}\\
  \textit{c08-modelo-io-01}\\

\end{flushright}
Como Haskell consegue realizar operações de entrada e saída e, ainda assim, preservar a ideia de linguagem puramente funcional?

\begin{enumerate}
  \item As funções de entrada e saída são exceções impuras dentro da linguagem e, por isso, não participam do sistema de tipos.

  \item Valores do tipo \haskellinline|IO a| são descrições de ações interativas. O programa constrói essas descrições de forma pura, e o \emph{runtime system} é responsável por executá-las.

  \item A notação \haskellinline|do| transforma temporariamente Haskell em uma linguagem imperativa, desativando a transparência referencial.

  \item O tipo \haskellinline|IO a| armazena internamente uma variável global mutável que pode ser alterada por qualquer função pura.

  \item Haskell não realiza entrada e saída; todo programa interativo precisa ser escrito em outra linguagem e apenas chamado a partir de Haskell.


\end{enumerate}

\questionTitle{10}{Inferência de tipo de isDigit}
\begin{flushright}
  \footnotesize
  \textit{\textbf{4: Tipos de Dados}}\\
  \textit{c04-inferencia-basica-02}\\

\end{flushright}
Considere \haskellinline{import Data.Char}. Qual é o tipo inferido para a expressão \haskellinline{isDigit '5'}?

\begin{enumerate}
  \item \haskellinline{Bool}

  \item \haskellinline{Char -> Bool}

  \item \haskellinline{Char}

  \item \haskellinline{True}

  \item \haskellinline{Num a => a}


\end{enumerate}

\questionTitle{11}{Estrutura geral de um laço recursivo de E/S}
\begin{flushright}
  \footnotesize
  \textit{\textbf{10: Ações de E/S recursivas}}\\
  \textit{c10-lacos-io-01}\\

\end{flushright}
No estilo usado no capítulo, qual é o padrão básico para repetir uma ação de E/S, como ler vários valores do teclado?

\begin{enumerate}
  \item Importar uma biblioteca obrigatória de laços, pois Haskell não permite repetição de ações de E/S sem ela.

  \item Declarar uma variável mutável global e atualizá-la até que a condição de parada seja falsa.

  \item Definir uma ação recursiva que executa uma parte da interação e chama a si mesma até atingir um caso base.

  \item Usar uma compreensão de listas dentro de \haskellinline|do| para repetir efeitos colaterais.

  \item Usar \haskellinline|map| sobre uma lista infinita e confiar que o \haskellinline|IO| executará apenas os elementos necessários.


\end{enumerate}

\questionTitle{12}{Regra para abreviação de comandos}
\begin{flushright}
  \footnotesize
  \textit{\textbf{2: Ambiente de desenvolvimento}}\\
  \textit{c02-ghci-inspecao-03}\\

\end{flushright}
No ambiente GHCi, os comandos que começam com dois-pontos (\texttt{:}) frequentemente podem ser abreviados (por exemplo, \texttt{:r} no lugar de \texttt{:reload}). Qual é a regra técnica exigida pelo ambiente para que uma abreviação seja aceita?

\begin{enumerate}
  \item O GHCi aceita qualquer abreviação independentemente de colisão; em caso de empate, ele sempre executa o último comando digitado no histórico da sessão.

  \item Apenas comandos que controlam o fluxo do terminal (como sair ou carregar) suportam atalhos; comandos de inspeção devem ser sempre digitados por extenso.

  \item A abreviação informada deve ser não ambígua, ou seja, deve ser um prefixo exclusivo que não coincida com a abreviação de nenhum outro comando existente.

  \item A abreviação só será aceita se o programador tiver configurado previamente um arquivo de manifesto indicando quais atalhos ele deseja mapear.

  \item A abreviação deve ter exatamente duas letras para que o analisador do GHCi não as confunda com nomes de variáveis Haskell válidas.


\end{enumerate}

\questionTitle{13}{Falha de readLn diante de entrada inválida}
\begin{flushright}
  \footnotesize
  \textit{\textbf{10: Ações de E/S recursivas}}\\
  \textit{c10-bordas-robustez-01}\\

\end{flushright}
No capítulo, os exemplos usam \haskellinline|readLn| para simplificar a leitura. O que acontece, em geral, se uma ação espera um número e o usuário digita um texto que não pode ser convertido para o tipo esperado?

\begin{enumerate}
  \item \haskellinline|readLn| repete a pergunta automaticamente até receber uma entrada válida.

  \item O valor inválido é interpretado como zero, pois zero é o valor neutro das leituras numéricas.

  \item O compilador detecta a entrada inválida antes da execução do programa.

  \item \haskellinline|readLn| devolve automaticamente \haskellinline|Nothing|, pois seu tipo padrão é \haskellinline|IO (Maybe a)|.

  \item A ação pode falhar em tempo de execução; uma versão mais robusta deveria validar a entrada, por exemplo com uma estratégia baseada em \haskellinline|readMaybe|.


\end{enumerate}

\questionTitle{14}{Comparação entre polimorfismo paramétrico e ad hoc}
\begin{flushright}
  \footnotesize
  \textit{\textbf{7: Polimorfismo}}\\
  \textit{cap07-comparacao-param-adhoc}\\

\end{flushright}
Qual alternativa distingue melhor polimorfismo paramétrico de polimorfismo \emph{ad hoc} em Haskell?
\begin{enumerate}
  \item Polimorfismo paramétrico é sinônimo de conversão implícita entre tipos numéricos.
  \item Polimorfismo \emph{ad hoc} não é verificado pelo sistema de tipos.
  \item No polimorfismo paramétrico, a função opera uniformemente para todos os tipos; no polimorfismo \emph{ad hoc}, o comportamento depende de uma instância de classe de tipo.
  \item No polimorfismo paramétrico, cada tipo precisa de uma implementação diferente; no \emph{ad hoc}, há sempre uma única implementação uniforme.
  \item Polimorfismo paramétrico ocorre apenas em POO; polimorfismo \emph{ad hoc} ocorre apenas em Haskell.

\end{enumerate}

\questionTitle{15}{Associatividade de subtração e potenciação}
\begin{flushright}
  \footnotesize
  \textit{\textbf{3: Expressões e definições}}\\
  \textit{c03-precedencia-associatividade-03}\\

\end{flushright}
Sabendo que \haskellinline|-| associa à esquerda e \haskellinline|^| associa à direita, qual alternativa está correta?

\begin{enumerate}
  \item \haskellinline|10 - 5 - 2| é lida como \haskellinline|(10 - 5) - 2|, enquanto \haskellinline|2 ^ 3 ^ 2| é lida como \haskellinline|2 ^ (3 ^ 2)|.

  \item Tanto \haskellinline|-| quanto \haskellinline|^| associam sempre à esquerda.

  \item A associatividade só se aplica a funções prefixas, não a operadores infixos.

  \item Tanto \haskellinline|-| quanto \haskellinline|^| associam sempre à direita.

  \item \haskellinline|10 - 5 - 2| é lida como \haskellinline|10 - (5 - 2)|, enquanto \haskellinline|2 ^ 3 ^ 2| é lida como \haskellinline|(2 ^ 3) ^ 2|.


\end{enumerate}

\questionTitle{16}{Comando para recarregar um arquivo no GHCi}
\begin{flushright}
  \footnotesize
  \textit{\textbf{3: Expressões e definições}}\\
  \textit{c03-ghci-definicoes-03}\\

\end{flushright}
Após carregar um arquivo Haskell no GHCi e editar o código no editor, qual comando recarrega o último arquivo carregado?

\begin{enumerate}
  \item \haskellinline|:load| sem argumento, pois ele sempre recarrega o arquivo anterior.

  \item \haskellinline|:browse|, pois ele atualiza o arquivo fonte em disco.

  \item \haskellinline|:type|, pois ele recompila e mostra os tipos de todas as definições.

  \item \haskellinline|:edit|, pois ele salva e recompila automaticamente o módulo.

  \item \haskellinline|:reload|, ou sua abreviação \haskellinline|:r|.


\end{enumerate}

\questionTitle{17}{Homogeneidade de listas}
\begin{flushright}
  \footnotesize
  \textit{\textbf{4: Tipos de Dados}}\\
  \textit{c04-checagem-listas-01}\\

\end{flushright}
Em Haskell, se um programador tentar avaliar a expressão \haskellinline{[1, True, "texto"]}, o que acontecerá?

\begin{enumerate}
  \item Ocorrerá um erro em tempo de compilação, pois todos os elementos de uma lista devem ter o mesmo tipo.

  \item A expressão será aceita e o tipo inferido será \haskellinline{Any}.

  \item Os valores serão implicitamente convertidos para \haskellinline{String}.

  \item Ocorrerá um erro apenas em tempo de execução, ao tentar usar os elementos.

  \item O compilador criará uma tupla automaticamente.


\end{enumerate}

\questionTitle{18}{Lista final produzida por acumulador com reverse}
\begin{flushright}
  \footnotesize
  \textit{\textbf{10: Ações de E/S recursivas}}\\
  \textit{c10-listas-ordem-03}\\

\end{flushright}
A função abaixo é executada inicialmente como \haskellinline|lerLista []|. O usuário digita \haskellinline|7, 1, 0|, um valor por linha.

\begin{haskellcode}
lerLista :: [Integer] -> IO [Integer]
lerLista xs = do
  x <- readLn
  if x == 0
    then return (reverse xs)
    else lerLista (x:xs)
\end{haskellcode}

Qual lista é produzida pela ação?

\begin{enumerate}
  \item \haskellinline|[1,7]|

  \item \haskellinline|[]|

  \item \haskellinline|[7,1,0]|

  \item \haskellinline|[7,1]|

  \item \haskellinline|[0]|


\end{enumerate}

\questionTitle{19}{Composição com fromMaybe}
\begin{flushright}
  \footnotesize
  \textit{\textbf{6: Estruturas de dados básicas}}\\
  \textit{c06-maybe-operacoes-01}\\

\end{flushright}
Qual é o resultado da expressão abaixo?
\begin{minted}{haskell}
fromMaybe "indefinido" (if True then Nothing else Just "ativo")
\end{minted}

\begin{enumerate}
  \item Ocorre um erro de tipos incompatíveis.

  \item \haskellinline{ Nothing }

  \item \haskellinline{ "ativo" }

  \item \haskellinline{ Just "ativo" }

  \item \haskellinline{ "indefinido" }


\end{enumerate}

\questionTitle{20}{Assinatura de função com Char, Int e String}
\begin{flushright}
  \footnotesize
  \textit{\textbf{4: Tipos de Dados}}\\
  \textit{c04-funcao-assinatura-03}\\

\end{flushright}
Qual é a assinatura de tipo de uma função que recebe um \haskellinline{Char}, um \haskellinline{Int} e retorna uma \haskellinline{String}?

\begin{enumerate}
  \item \haskellinline{(Char, Int) -> String}

  \item \haskellinline{[Char, Int, String]}

  \item \haskellinline{String -> Char -> Int}

  \item \haskellinline{Char -> Int -> String}

  \item \haskellinline{Char, Int -> String}


\end{enumerate}

\questionTitle{21}{Quando usar realToFrac}
\begin{flushright}
  \footnotesize
  \textit{\textbf{7: Polimorfismo}}\\
  \textit{cap07-conversao-realToFrac}\\

\end{flushright}
Qual é o papel de \haskellinline{realToFrac} em conversões numéricas?
\begin{enumerate}
  \item Converter apenas de \haskellinline{Int} para \haskellinline{Integer}, preservando integralidade.
  \item Permitir que operadores como \haskellinline{(+)} misturem automaticamente \haskellinline{Int} e \haskellinline{Double}.
  \item Converter de um tipo instância de \haskellinline{Real} para um tipo de destino instância de \haskellinline{Fractional}, como de \haskellinline{Float} para \haskellinline{Double}.
  \item Arredondar qualquer número fracionário para o inteiro mais próximo.
  \item Substituir \haskellinline{read} na conversão de \haskellinline{String} para número.

\end{enumerate}

\questionTitle{22}{Explorando módulos com :browse}
\begin{flushright}
  \footnotesize
  \textit{\textbf{2: Ambiente de desenvolvimento}}\\
  \textit{c02-ghci-modulos-03}\\

\end{flushright}
Após adicionar um módulo à sessão atual do GHCi, o desenvolvedor digita o comando \texttt{:browse Data.List}. Qual será o resultado visível no terminal?

\begin{enumerate}
  \item O GHCi listará as assinaturas de tipo de todas as funções e definições de dados que estão exportadas e disponíveis publicamente naquele módulo.

  \item O GHCi exibirá o código-fonte C gerado nos bastidores para cada função presente em \texttt{Data.List}.

  \item O GHCi procurará erros de compilação (\emph{warnings}) ignorados nas funções internas daquele módulo específico.

  \item O GHCi removerá o módulo do escopo ativo, fechando as portas para futuras avaliações matemáticas.

  \item O GHCi abrirá uma nova aba no navegador web (\emph{browser}) do sistema operacional exibindo a documentação HTML oficial do módulo.


\end{enumerate}

\questionTitle{23}{Prompt sem quebra de linha}
\begin{flushright}
  \footnotesize
  \textit{\textbf{8: Programas interativos}}\\
  \textit{c08-buffering-01}\\

\end{flushright}
Um programa usa \haskellinline|putStr| para imprimir um \emph{prompt} e, em seguida, chama \haskellinline|getLine|. Em alguns ambientes, o \emph{prompt} pode não aparecer antes da digitação. Qual é uma solução adequada?

\begin{enumerate}
  \item Forçar a descarga do \emph{buffer} com \haskellinline|hFlush stdout| logo após \haskellinline|putStr|, ou configurar \haskellinline|stdout| com \haskellinline|NoBuffering|.

  \item Trocar \haskellinline|getLine| por \haskellinline|readLn|, pois \haskellinline|readLn| sempre descarrega a saída padrão.

  \item Adicionar \haskellinline|let| antes de \haskellinline|getLine| para forçar a ordem de execução.

  \item Reordenar o programa colocando \haskellinline|getLine| antes de \haskellinline|putStr|.

  \item Trocar \haskellinline|putStr| por \haskellinline|print|, pois \haskellinline|print| não usa \emph{buffer}.


\end{enumerate}

\questionTitle{24}{Trabalho pendente após a chamada recursiva}
\begin{flushright}
  \footnotesize
  \textit{\textbf{10: Ações de E/S recursivas}}\\
  \textit{c10-acumuladores-cauda-01}\\

\end{flushright}
Considere a versão sem acumulador:

\begin{haskellcode}
lerESomar :: IO Integer
lerESomar = do
  n <- readLn
  if n == 0
    then return 0
    else do
      somaResto <- lerESomar
      return (n + somaResto)
\end{haskellcode}

Por que a chamada recursiva não está em posição de cauda?

\begin{enumerate}
  \item Porque qualquer uso de \haskellinline|readLn| impede uma função de ser recursiva de cauda.

  \item Porque \haskellinline|if| não permite chamadas recursivas em posição de cauda.

  \item Porque o tipo \haskellinline|IO Integer| é incompatível com recursão de cauda.

  \item Porque o caso base usa \haskellinline|return 0| em vez de \haskellinline|return n|.

  \item Porque, depois que \haskellinline|lerESomar| retorna, ainda é necessário calcular \haskellinline|n + somaResto| e aplicar \haskellinline|return| ao resultado.


\end{enumerate}

\questionTitle{25}{Inferência de restrições implícitas}
\begin{flushright}
  \footnotesize
  \textit{\textbf{7: Polimorfismo}}\\
  \textit{cap07-inferencia-classes-implicitas}\\

\end{flushright}
Considere a função:

\begin{minted}{haskell}
positivoMaiorQue x y = x > y && y > 0
\end{minted}

Qual é a assinatura mais geral?

\begin{enumerate}
  \item \mintinline{haskell}{positivoMaiorQue :: Ord a => a -> a -> Bool}
  \item \mintinline{haskell}{positivoMaiorQue :: Integral a => a -> a -> Bool}
  \item \mintinline{haskell}{positivoMaiorQue :: Num a => a -> a -> Bool}
  \item \mintinline{haskell}{positivoMaiorQue :: Bool -> Bool -> Bool}
  \item \mintinline{haskell}{positivoMaiorQue :: (Ord a, Num a) => a -> a -> Bool}

\end{enumerate}

\questionTitle{26}{Tipo da condição}
\begin{flushright}
  \footnotesize
  \textit{\textbf{5: Expressão condicional}}\\
  \textit{c05-if-regras-02}\\

\end{flushright}
Qual das expressões abaixo falha especificamente porque a condição do \haskellinline|if| não tem tipo \haskellinline|Bool|?

\begin{enumerate}
  \item \haskellinline|if even 4 then 10 else 20|

  \item \haskellinline|if 3 then 10 else 20|

  \item \haskellinline|if 3 < 5 then 10 else 20|

  \item \haskellinline|if True then 10 else 20|

  \item \haskellinline|if null [] then 10 else 20|


\end{enumerate}

\questionTitle{27}{Valores e forma normal}
\begin{flushright}
  \footnotesize
  \textit{\textbf{3: Expressões e definições}}\\
  \textit{c03-comentarios-valores-03}\\

\end{flushright}
No contexto da avaliação de expressões em Haskell, qual alternativa caracteriza corretamente a relação entre expressões e valores?

\begin{enumerate}
  \item Uma expressão só é considerada valor se tiver sido associada a um nome por uma definição.

  \item Todo valor é uma expressão, mas nem toda expressão já é um valor; uma expressão torna-se um valor quando atinge uma forma que não pode mais ser reduzida.

  \item Uma expressão em forma normal ainda precisa ser reduzida pelo menos uma vez para produzir um valor.

  \item Valores são apenas números; listas, tuplas e funções não são considerados valores em Haskell.

  \item Toda expressão escrita pelo programador já é imediatamente um valor, mesmo antes de qualquer avaliação.


\end{enumerate}

\questionTitle{28}{Conversão de operador infixo para função prefixa}
\begin{flushright}
  \footnotesize
  \textit{\textbf{3: Expressões e definições}}\\
  \textit{c03-aplicacao-notacoes-03}\\

\end{flushright}
Qual alternativa é uma reescrita prefixa equivalente à expressão \haskellinline|100 `mod` (5 * 2)|?

\begin{enumerate}
  \item \haskellinline|100 mod (5 * 2)|.

  \item \haskellinline|mod (100, 5 * 2)|.

  \item \haskellinline|(mod 100 5) * 2|.

  \item \haskellinline|mod (100 * 5) 2|.

  \item \haskellinline|mod 100 (5 * 2)|.


\end{enumerate}

\questionTitle{29}{Acesso composto a pares}
\begin{flushright}
  \footnotesize
  \textit{\textbf{6: Estruturas de dados básicas}}\\
  \textit{c06-tuplas-acesso-01}\\

\end{flushright}
Considere a expressão abaixo.
\begin{minted}{haskell}
snd (fst (("Ana", 21), True))
\end{minted}
Qual é o seu resultado?

\begin{enumerate}
  \item \haskellinline{ 21 }

  \item A expressão resulta em erro de compilação.

  \item \haskellinline{ "Ana" }

  \item \haskellinline{ ("Ana",21) }

  \item \haskellinline{ True }


\end{enumerate}

\questionTitle{30}{Restrições combinadas em função sobrecarregada}
\begin{flushright}
  \footnotesize
  \textit{\textbf{7: Polimorfismo}}\\
  \textit{cap07-adhoc-funcao-inferida-num-show}\\

\end{flushright}
Considere a definição:

\begin{minted}{haskell}
f x = show (x + 1)
\end{minted}

Qual assinatura de tipo é a mais geral?

\begin{enumerate}
  \item \mintinline{haskell}{f :: (Num a, Show a) => a -> String}
  \item \mintinline{haskell}{f :: Show a => a -> String}
  \item \mintinline{haskell}{f :: a -> String}
  \item \mintinline{haskell}{f :: Num a => a -> String}
  \item \mintinline{haskell}{f :: Int -> Int}

\end{enumerate}

\questionTitle{31}{Ordem das definições locais em where}
\begin{flushright}
  \footnotesize
  \textit{\textbf{3: Expressões e definições}}\\
  \textit{c03-definicoes-where-03}\\

\end{flushright}
Considere a definição:

\begin{haskellcode}
g x = a * b
  where
    b = x + 2
    a = b - 1
\end{haskellcode}

Qual é o valor de \haskellinline|g 4|?

\begin{enumerate}
  \item \haskellinline|20|.

  \item \haskellinline|18|.

  \item Ocorre erro, pois \haskellinline|a| aparece antes de \haskellinline|b| no corpo da função.

  \item Ocorre erro, pois definições em \haskellinline|where| não podem depender umas das outras.

  \item \haskellinline|30|.


\end{enumerate}

\questionTitle{32}{Rastreamento do fatorial}
\begin{flushright}
  \footnotesize
  \textit{\textbf{9: Funções recursivas}}\\
  \textit{c09-rastreamento-fatorial-01}\\

\end{flushright}
Considere:
\begin{haskellcode}
fatorial :: Integer -> Integer
fatorial n
  | n == 0 = 1
  | n > 0  = n * fatorial (n - 1)
\end{haskellcode}
Qual é o valor de \haskellinline|fatorial 4|?

\begin{enumerate}
  \item \haskellinline|12|

  \item \haskellinline|20|

  \item \haskellinline|16|

  \item \haskellinline|24|

  \item \haskellinline|120|


\end{enumerate}

\questionTitle{33}{Literais fracionários e fromRational}
\begin{flushright}
  \footnotesize
  \textit{\textbf{7: Polimorfismo}}\\
  \textit{cap07-literais-fromrational-contexto}\\

\end{flushright}
Qual afirmação descreve corretamente literais fracionários como \haskellinline{3.14}?
\begin{enumerate}
  \item Eles podem ser usados diretamente como \haskellinline{Int}, pois Haskell arredonda automaticamente.
  \item Eles exigem \haskellinline{Integral a}, pois possuem uma representação decimal.
  \item Eles são convertidos por \haskellinline{fromInteger}, assim como \haskellinline{42}.
  \item Eles são interpretados por meio de \haskellinline{fromRational} e recebem inicialmente uma restrição \haskellinline{Fractional a => a}.
  \item Eles têm sempre tipo \haskellinline{Double}, independentemente do contexto.

\end{enumerate}

\questionTitle{34}{Equivalência com True}
\begin{flushright}
  \footnotesize
  \textit{\textbf{5: Expressão condicional}}\\
  \textit{c05-otherwise-02}\\

\end{flushright}
Qual definição é equivalente ao uso usual de \haskellinline|otherwise| na última guarda?

\begin{enumerate}
  \item Substituir \haskellinline|otherwise| por \haskellinline|True|.

  \item Substituir \haskellinline|otherwise| por \haskellinline|else|.

  \item Substituir \haskellinline|otherwise| por \haskellinline|False|.

  \item Substituir \haskellinline|otherwise| por \haskellinline|undefined|.

  \item Substituir \haskellinline|otherwise| por \haskellinline|()|.


\end{enumerate}

\questionTitle{35}{Inferência com zip, head e fst}
\begin{flushright}
  \footnotesize
  \textit{\textbf{7: Polimorfismo}}\\
  \textit{cap07-param-zip-fst-inferencia}\\

\end{flushright}
Considere:

\begin{minted}{haskell}
f xs ys = fst (head (zip xs ys))
\end{minted}

Qual é a assinatura de tipo mais geral de \haskellinline{f}?

\begin{enumerate}
  \item \mintinline{haskell}{f :: [(a,b)] -> a}
  \item \mintinline{haskell}{f :: [a] -> [a] -> a}
  \item \mintinline{haskell}{f :: Ord a => [a] -> [b] -> a}
  \item \mintinline{haskell}{f :: [a] -> [b] -> a}
  \item \mintinline{haskell}{f :: [a] -> [b] -> b}

\end{enumerate}

\questionTitle{36}{let e <- em blocos do}
\begin{flushright}
  \footnotesize
  \textit{\textbf{8: Programas interativos}}\\
  \textit{c08-blocos-do-01}\\

\end{flushright}
Dentro de um bloco \haskellinline|do| de uma ação \haskellinline|IO|, qual é a diferença fundamental entre \haskellinline|x <- acao| e \haskellinline|let x = expr|?

\begin{enumerate}
  \item \haskellinline|x <- acao| executa uma ação de E/S e associa a \haskellinline|x| o valor produzido por ela; \haskellinline|let x = expr| define localmente um valor ou função pura, sem executar ação de E/S.

  \item \haskellinline|let| só é permitido fora de blocos \haskellinline|do|.

  \item \haskellinline|let| executa a ação de E/S, enquanto \haskellinline|<-| apenas cria um apelido para a ação.

  \item Não há diferença semântica; as duas formas são intercambiáveis em blocos \haskellinline|do|.

  \item \haskellinline|<-| só pode ser usado para valores numéricos; \haskellinline|let| só pode ser usado para strings.


\end{enumerate}

\questionTitle{37}{Análise declarativa da abstração do Quicksort}
\begin{flushright}
  \footnotesize
  \textit{\textbf{1: Paradigmas de programação}}\\
  \textit{c01-mecanismo-reducao-02}\\

\end{flushright}
Considere a cláusula principal da implementação declarativa do \emph{Quicksort} em Haskell:

\begin{haskellcode}
qsort (pivo:resto) = qsort (filter (<= pivo) resto) ++ 
                     [pivo] ++ 
                     qsort (filter (> pivo) resto)
\end{haskellcode}

Ao comparar este trecho com a tradicional implementação iterativa em C, qual das seguintes afirmações avalia corretamente a diferença de abstração utilizada na partição dos elementos?

\begin{enumerate}
  \item A versão funcional requer que a lista de entrada seja previamente ordenada, caso contrário a função \haskellinline|filter| lançará uma exceção de limite de tempo de execução.

  \item A versão funcional abstrai a manipulação de índices de memória e laços \cinline|while|, descrevendo a lista particionada como o resultado da aplicação de funções de filtragem baseadas em propriedades lógicas da lista.

  \item A versão funcional oculta a passagem de ponteiros de memória nos bastidores do operador \haskellinline|:|, realizando uma mutação de array idêntica à troca de variáveis (\cinline|swap|) do algoritmo em C.

  \item A versão funcional é inerentemente mais ineficiente pois exige que o compilador traduza o operador de concatenação (\haskellinline|++|) em loops duplamente encadeados estritos.

  \item A versão funcional falha em implementar a estratégia de \emph{divisão e conquista}, pois a linguagem Haskell não suporta chamadas de recursão de cauda nativa como a linguagem C.


\end{enumerate}

\questionTitle{38}{Strings como listas}
\begin{flushright}
  \footnotesize
  \textit{\textbf{6: Estruturas de dados básicas}}\\
  \textit{c06-listas-operacoes-02}\\

\end{flushright}
Como \haskellinline{String} é uma lista de caracteres, a expressão abaixo pode ser avaliada pelas mesmas funções de listas.
\begin{minted}{haskell}
drop 5 "Paradigma"
\end{minted}
Qual é o resultado?

\begin{enumerate}
  \item A expressão resulta em erro de tipo.

  \item \haskellinline{ "Parad" }

  \item \haskellinline{ "gma" }

  \item \haskellinline{ "aradigma" }

  \item \haskellinline{ "igma" }


\end{enumerate}

\questionTitle{39}{Identificadores alfanuméricos válidos}
\begin{flushright}
  \footnotesize
  \textit{\textbf{3: Expressões e definições}}\\
  \textit{c03-arquivos-identificadores-03}\\

\end{flushright}
De acordo com as regras para identificadores alfanuméricos usados como nomes de variáveis ou funções, qual alternativa apresenta um identificador válido?

\begin{enumerate}
  \item \haskellinline|ValorFinal|.

  \item \haskellinline|2valor|.

  \item \haskellinline|module|.

  \item \haskellinline|valor_final'|.

  \item \haskellinline|valor-final|.


\end{enumerate}

\questionTitle{40}{Otimização de chamada de cauda em Haskell}
\begin{flushright}
  \footnotesize
  \textit{\textbf{9: Funções recursivas}}\\
  \textit{c09-eficiencia-tco-01}\\

\end{flushright}
Qual afirmação é a mais adequada sobre recursividade de cauda e uso de memória em Haskell?

\begin{enumerate}
  \item Recursividade de cauda é impossível em Haskell, pois a linguagem é puramente funcional.

  \item Toda função recursiva de cauda em Haskell usa memória constante, independentemente da avaliação dos argumentos.

  \item O uso de acumuladores sempre elimina a necessidade de casos base.

  \item Haskell rejeita qualquer função que não seja recursiva de cauda.

  \item Recursividade de cauda ajuda a evitar operações pendentes, mas em Haskell o uso de memória também depende de como os acumuladores são avaliados.


\end{enumerate}

\questionTitle{41}{Cabeça, cauda e construção}
\begin{flushright}
  \footnotesize
  \textit{\textbf{6: Estruturas de dados básicas}}\\
  \textit{c06-listas-estrutura-02}\\

\end{flushright}
Qual é o resultado da expressão abaixo?
\begin{minted}{haskell}
head ([10,20] ++ [30,40])
\end{minted}

\begin{enumerate}
  \item \haskellinline{ 20 }

  \item \haskellinline{ 30 }

  \item \haskellinline{ 10 }

  \item A expressão resulta em erro de execução.

  \item \haskellinline{ [10,20] }


\end{enumerate}

\questionTitle{42}{String como lista de Char}
\begin{flushright}
  \footnotesize
  \textit{\textbf{4: Tipos de Dados}}\\
  \textit{c04-tipos-string-01}\\

\end{flushright}
Sabendo que \haskellinline{String} é um tipo sinônimo para \haskellinline{[Char]}, qual das seguintes expressões é \textbf{equivalente} a \haskellinline{"FP"}?

\begin{enumerate}
  \item \haskellinline{'FP'}

  \item \haskellinline{"F" : "P"}

  \item \haskellinline{['F', 'P']}

  \item \haskellinline{('F', 'P')}

  \item \haskellinline{["F", "P"]}


\end{enumerate}

\questionTitle{43}{A ação readLn}
\begin{flushright}
  \footnotesize
  \textit{\textbf{8: Programas interativos}}\\
  \textit{c08-conversao-validacao-01}\\

\end{flushright}
O que a ação \haskellinline|readLn :: Read a => IO a| faz quando é executada?

\begin{enumerate}
  \item Ela apenas lê uma linha de texto, sendo idêntica a \haskellinline|getLine|.

  \item Ela lê uma linha da entrada padrão e tenta convertê-la para um valor de algum tipo pertencente à classe \haskellinline|Read|.

  \item Ela valida a entrada automaticamente e, em caso de erro, pede uma nova tentativa.

  \item Ela lê um único caractere da entrada padrão.

  \item Ela imprime a representação textual de um valor, como \haskellinline|print|.


\end{enumerate}

\questionTitle{44}{A Crise do Software e a solução funcional}
\begin{flushright}
  \footnotesize
  \textit{\textbf{1: Paradigmas de programação}}\\
  \textit{c01-engenharia-software-02}\\

\end{flushright}
O termo \emph{Crise do Software} surgiu historicamente devido à dificuldade de entregar sistemas complexos no prazo, dentro do orçamento e isentos de falhas catastróficas em tempo de execução. Como o paradigma funcional atua para mitigar diretamente a raiz estrutural dessa crise?

\begin{enumerate}
  \item Restringindo rigorosamente a arquitetura dos projetos a um limite máximo de linhas de código por pacote para evitar que os softwares atinjam uma escala insustentável.

  \item Reduzindo o tempo de desenvolvimento ao automatizar a geração de código através de modelos probabilísticos baseados em inteligência artificial no sistema de tipos dinâmicos.

  \item Forçando a adoção de metodologias de desenvolvimento ágil (\emph{Agile/Scrum}) através de ferramentas nativas embutidas diretamente nos compiladores das linguagens funcionais.

  \item Eliminando de forma definitiva a necessidade de escrever casos de teste unitário, já que o processo de compilação cruzada garante a integridade incondicional das regras de negócio.

  \item Elevando o nível de abstração focado na declaração (o que fazer) e adotando bases que eliminam o estado mutável não documentado, o que reduz drasticamente a complexidade do rastreamento de defeitos.


\end{enumerate}

\questionTitle{45}{Interpretação de Num a => a}
\begin{flushright}
  \footnotesize
  \textit{\textbf{4: Tipos de Dados}}\\
  \textit{c04-ghci-interpretacao-01}\\

\end{flushright}
Ao executar \texttt{:type 42} no GHCi, a resposta é \haskellinline{42 :: Num a => a}. O que isso significa?

\begin{enumerate}
  \item \haskellinline{42} é uma constante numérica que só pode ser do tipo \haskellinline{Int}.

  \item Significa que \haskellinline{42} é uma função que recebe um \haskellinline{Num} e retorna um \haskellinline{a}.

  \item \haskellinline{42} pode ser interpretado como um valor de qualquer tipo que pertença à classe \haskellinline{Num} (como \haskellinline{Int}, \haskellinline{Integer}, \haskellinline{Double}).

  \item Indica um erro de tipo, pois o tipo é ambíguo e não pode ser resolvido.

  \item \haskellinline{42} é um valor de um tipo desconhecido que não pode ser usado em operações aritméticas.


\end{enumerate}

\questionTitle{46}{O significado de IO ()}
\begin{flushright}
  \footnotesize
  \textit{\textbf{8: Programas interativos}}\\
  \textit{c08-tipos-acoes-01}\\

\end{flushright}
Por que muitas ações de saída, como \haskellinline|putStrLn :: String -> IO ()|, produzem um valor do tipo \haskellinline|IO ()|?

\begin{enumerate}
  \item Porque a ação realiza um efeito de E/S, mas não produz nenhum dado significativo para ser usado pelo restante do programa. O valor produzido é o valor unitário \haskellinline|()|.

  \item Porque o tipo \haskellinline|()| indica que a ação pode ser executada no máximo uma vez.

  \item Porque a função não recebe argumentos.

  \item Porque a ação sempre falha e retorna um valor vazio.

  \item Porque \haskellinline|IO ()| é um sinônimo de \haskellinline|IO String|.


\end{enumerate}

\questionTitle{47}{Lookup com valor padrão}
\begin{flushright}
  \footnotesize
  \textit{\textbf{6: Estruturas de dados básicas}}\\
  \textit{c06-lookup-associacao-02}\\

\end{flushright}
Qual é o resultado da expressão abaixo?
\begin{minted}{haskell}
let tabela = [("ana",6.5),("bia",8.0)]
in fromMaybe 0.0 (lookup "paula" tabela)
\end{minted}

\begin{enumerate}
  \item \haskellinline{ 8.0 }

  \item \haskellinline{ Just 0.0 }

  \item \haskellinline{ Nothing }

  \item Ocorre um erro em tempo de execução.

  \item \haskellinline{ 0.0 }


\end{enumerate}

\questionTitle{48}{Modelagem de inventário}
\begin{flushright}
  \footnotesize
  \textit{\textbf{6: Estruturas de dados básicas}}\\
  \textit{c06-modelagem-02}\\

\end{flushright}
Você quer representar um inventário como uma lista de produtos. Cada produto possui nome (\haskellinline{String}), código (\haskellinline{Int}) e preço opcional (\haskellinline{Maybe Double}). Qual tipo representa corretamente esse inventário?

\begin{enumerate}
  \item \haskellinline{Maybe (String, Int, Double)}

  \item \haskellinline{[(String, Int, Maybe Double)]}

  \item \haskellinline{(String, Int, Maybe Double)}

  \item \haskellinline{([String], [Int], [Maybe Double])}

  \item \haskellinline{[String, Int, Maybe Double]}


\end{enumerate}

\questionTitle{49}{Refatoração para acumulador}
\begin{flushright}
  \footnotesize
  \textit{\textbf{9: Funções recursivas}}\\
  \textit{c09-cauda-refatoracao-01}\\

\end{flushright}
Qual alternativa é uma refatoração correta de \haskellinline|fatorial| usando uma função auxiliar com acumulador?

\begin{enumerate}
  \item \begin{haskellcode}
fat n = fat' n 1
  where
    fat' n acc
      | n == 0 = acc
      | n > 0  = fat' (n - 1) (n * acc)
\end{haskellcode}

  \item \begin{haskellcode}
fat n = fat' n 1
  where
    fat' n acc
      | n == 0 = 1
      | n > 0  = fat' (n + 1) (n * acc)
\end{haskellcode}

  \item \begin{haskellcode}
fat n = fat' n 1
  where
    fat' n acc
      | n == 0 = n
      | n > 0  = fat' (n - 1) acc
\end{haskellcode}

  \item \begin{haskellcode}
fat n = fat' n 0
  where
    fat' n acc
      | n == 0 = acc
      | n > 0  = fat' (n - 1) (n * acc)
\end{haskellcode}

  \item \begin{haskellcode}
fat n = fat' n 1
  where
    fat' n acc
      | n == 0 = acc
      | n > 0  = n * fat' (n - 1) acc
\end{haskellcode}


\end{enumerate}

\questionTitle{50}{Ordem de avaliação das guardas}
\begin{flushright}
  \footnotesize
  \textit{\textbf{5: Expressão condicional}}\\
  \textit{c05-guardas-sintaxe-semantica-02}\\

\end{flushright}
Como as guardas de uma definição são avaliadas?

\begin{enumerate}
  \item O compilador reordena as guardas para testar primeiro as mais específicas.

  \item As guardas são avaliadas de baixo para cima.

  \item Todas as guardas são avaliadas, e a última verdadeira determina o resultado.

  \item De cima para baixo; a primeira guarda verdadeira determina o resultado.

  \item O resultado é escolhido pela guarda cuja expressão à direita de \haskellinline|=| parece mais simples.


\end{enumerate}

\questionTitle{51}{A analogia da receita de bolo}
\begin{flushright}
  \footnotesize
  \textit{\textbf{1: Paradigmas de programação}}\\
  \textit{c01-conceitos-paradigmas-03}\\

\end{flushright}
Ao explicar paradigmas de programação, é comum utilizar a analogia de uma \emph{receita de bolo} (ex: \emph{1. quebre os ovos; 2. adicione farinha; 3. asse por 40 minutos}). 

Esta analogia ilustra de forma primária o raciocínio de qual paradigma e por quê?

\begin{enumerate}
  \item O paradigma declarativo, pois descreve as propriedades finais do bolo sem se preocupar com a ordem matemática em que os ingredientes foram misturados.

  \item O paradigma orientado a objetos, pois a receita representa a troca de mensagens assíncronas entre as instâncias dos ingredientes.

  \item O paradigma lógico, pois cada passo da receita atua como um axioma em um motor de inferência que deduz a existência do bolo.

  \item O paradigma imperativo, pois foca na descrição passo a passo da sequência de controle e nas mudanças sucessivas de estado para atingir o objetivo.

  \item O paradigma funcional puro, pois cada passo da receita é uma função referencialmente transparente que não altera o estado da cozinha.


\end{enumerate}

\questionTitle{52}{Refatoração para forma mais concisa}
\begin{flushright}
  \footnotesize
  \textit{\textbf{5: Expressão condicional}}\\
  \textit{c05-if-vs-guardas-02}\\

\end{flushright}
Considere a função:
\begin{haskellcode}
classifica n
  | n > 0     = "Positivo"
  | n < 0     = "Negativo"
  | n == 0    = "Zero"
  | otherwise = "Impossível"
\end{haskellcode}
Qual definição é funcionalmente equivalente, mas mais concisa?

\begin{enumerate}
  \item \begin{haskellcode}
classifica n
  | n > 0     = "Positivo"
  | otherwise = "Negativo"
\end{haskellcode}

  \item \begin{haskellcode}
classifica n
  | n > 0     = "Positivo"
  | n < 0     = "Negativo"
  | otherwise = "Zero"
\end{haskellcode}

  \item A função original não pode ser simplificada sem mudar seu comportamento.

  \item \begin{haskellcode}
classifica n
  | n == 0    = "Zero"
  | otherwise = "Positivo ou Negativo"
\end{haskellcode}

  \item \begin{haskellcode}
classifica n
  | n >= 0    = "Positivo"
  | otherwise = "Negativo"
\end{haskellcode}


\end{enumerate}

\questionTitle{53}{Segurança de tipos com sinônimos}
\begin{flushright}
  \footnotesize
  \textit{\textbf{4: Tipos de Dados}}\\
  \textit{c04-sinonimos-seguranca-01}\\

\end{flushright}
Considere as definições \haskellinline{type UsuarioID = Int} e \haskellinline{type ProdutoID = Int}. Se uma função espera um \haskellinline{UsuarioID} mas recebe um \haskellinline{ProdutoID}, o que acontecerá?

\begin{enumerate}
  \item A palavra-chave \haskellinline{type} não pode ser usada com tipos numéricos.

  \item O programa falhará na compilação com um erro de incompatibilidade de tipos.

  \item O programa compilará normalmente, pois para o sistema de tipos, \haskellinline{UsuarioID}, \haskellinline{ProdutoID} e \haskellinline{Int} são o mesmo tipo.

  \item O programa compilará, mas emitirá um aviso (\emph{warning}).

  \item Ocorrerá um erro em tempo de execução.


\end{enumerate}

\questionTitle{54}{Saudação interativa}
\begin{flushright}
  \footnotesize
  \textit{\textbf{8: Programas interativos}}\\
  \textit{c08-programas-completos-01}\\

\end{flushright}
Considere o programa:

\begin{minted}{haskell}
main :: IO ()
main = do
  putStrLn "Nome?"
  nome <- getLine
  putStrLn ("Olá, " ++ nome ++ "!")
\end{minted}

Se o usuário digitar \haskellinline|Ana|, qual será a última linha impressa?

\begin{enumerate}
  \item \haskellinline|Ana|

  \item \haskellinline|Nome?|

  \item \haskellinline|Olá, nome!|

  \item \haskellinline|Olá, getLine!|

  \item \haskellinline|Olá, Ana!|


\end{enumerate}

\questionTitle{55}{Let como subexpressão}
\begin{flushright}
  \footnotesize
  \textit{\textbf{3: Expressões e definições}}\\
  \textit{c03-let-escopo-03}\\

\end{flushright}
Qual é o valor da expressão \haskellinline|2 * (let x = 3; y = 4 in x + y)|?

\begin{enumerate}
  \item A expressão é inválida, pois \haskellinline|let| não pode aparecer dentro de parênteses.

  \item \haskellinline|14|.

  \item \haskellinline|24|.

  \item \haskellinline|9|.

  \item \haskellinline|7|.


\end{enumerate}

\questionTitle{56}{Primeira guarda verdadeira vence}
\begin{flushright}
  \footnotesize
  \textit{\textbf{5: Expressão condicional}}\\
  \textit{c05-guardas-avaliacao-02}\\

\end{flushright}
Considere a função:
\begin{haskellcode}
f x
  | x >= 0    = 10
  | x > 3     = 20
  | otherwise = 30
\end{haskellcode}
Qual é o valor de \haskellinline|f 4|?

\begin{enumerate}
  \item \haskellinline|20|

  \item O compilador rejeita a definição pela sobreposição das guardas.

  \item \haskellinline|10|

  \item Ocorre um erro, pois duas guardas são verdadeiras.

  \item \haskellinline|30|


\end{enumerate}

\questionTitle{57}{Avaliação com cálculo local}
\begin{flushright}
  \footnotesize
  \textit{\textbf{5: Expressão condicional}}\\
  \textit{c05-where-avaliacao-02}\\

\end{flushright}
Considere a função:
\begin{haskellcode}
analisa y
  | x > 10    = "Alto"
  | otherwise = "Baixo"
  where
    x = y * y + 1
\end{haskellcode}
Qual é o resultado de \haskellinline|analisa 3|?

\begin{enumerate}
  \item Ocorre um erro de tipo em \haskellinline|x > 10|.

  \item \haskellinline|"Baixo"|

  \item \haskellinline|"Alto"|

  \item \haskellinline|10|

  \item \haskellinline|3|


\end{enumerate}

\questionTitle{58}{Ineficiência do Fibonacci ingênuo}
\begin{flushright}
  \footnotesize
  \textit{\textbf{9: Funções recursivas}}\\
  \textit{c09-formas-fibonacci-01}\\

\end{flushright}
Considere a definição ingênua:
\begin{haskellcode}
fib n | n == 0    = 0
      | n == 1    = 1
      | otherwise = fib (n - 2) + fib (n - 1)
\end{haskellcode}
Por que essa definição é ineficiente para valores grandes de \haskellinline|n|?

\begin{enumerate}
  \item Porque ela recalcula os mesmos subproblemas muitas vezes, fazendo o número de chamadas crescer rapidamente.

  \item Porque \haskellinline|+| entre inteiros é proibido em funções recursivas.

  \item Porque a chamada recursiva é de cauda e, portanto, nunca retorna.

  \item Porque ela armazena todos os resultados anteriores em uma lista global.

  \item Porque a função não tem casos base para \haskellinline|0| e \haskellinline|1|.


\end{enumerate}

\questionTitle{59}{Projeto para cálculo de média}
\begin{flushright}
  \footnotesize
  \textit{\textbf{10: Ações de E/S recursivas}}\\
  \textit{c10-separacao-puro-03}\\

\end{flushright}
Para o exercício da média de uma sequência, considere o objetivo: ler valores até um sentinela negativo e imprimir a média dos valores não negativos lidos. Qual decomposição é mais alinhada ao estilo do capítulo?

\begin{enumerate}
  \item Uma função \haskellinline|sum :: IO [Double] -> Double|, pois \haskellinline|sum| opera diretamente sobre ações de E/S.

  \item Uma ação \haskellinline|main| que acumula uma string com todos os números e chama \haskellinline|read| apenas no final.

  \item Uma ação \haskellinline|lerLista :: IO [Double]| para coletar os valores e uma função pura ou expressão separada para calcular a média da lista, tratando a lista vazia antes da divisão.

  \item Uma única função pura \haskellinline|media :: IO Double| que chama \haskellinline|readLn| internamente sem usar \haskellinline|do|.

  \item Uma definição que ignora a lista vazia, pois \haskellinline|fromIntegral (length lista)| nunca pode ser zero.


\end{enumerate}

\questionTitle{60}{Caso base inalcançável}
\begin{flushright}
  \footnotesize
  \textit{\textbf{9: Funções recursivas}}\\
  \textit{c09-fundamentos-terminacao-01}\\

\end{flushright}
Considere a definição abaixo:
\begin{haskellcode}
fatorial :: Integer -> Integer
fatorial n
  | n == 0 = 1
  | otherwise = n * fatorial (n - 1)
\end{haskellcode}
O que acontece ao avaliar \haskellinline|fatorial (-1)|?

\begin{enumerate}
  \item A computação não alcança o caso base: ela chama \haskellinline|fatorial (-2)|, depois \haskellinline|fatorial (-3)|, e assim por diante.

  \item A função retorna \haskellinline|0|, pois fatoriais de números negativos são tratados como zero em Haskell.

  \item A função retorna \haskellinline|1|, pois \haskellinline|0| é o caso base mais próximo.

  \item Ocorre imediatamente um erro de padrões não exaustivos, antes de qualquer chamada recursiva.

  \item A função calcula \haskellinline|fatorial 1|, ignorando automaticamente o sinal negativo.


\end{enumerate}

\questionTitle{61}{Análise de restrições combinadas}
\begin{flushright}
  \footnotesize
  \textit{\textbf{7: Polimorfismo}}\\
  \textit{cap07-expressao-comparacao-show}\\

\end{flushright}
Considere:

\begin{minted}{haskell}
h x y = x <= y && show x == show y
\end{minted}

Qual é a assinatura mais geral de \haskellinline{h}?

\begin{enumerate}
  \item \mintinline{haskell}{h :: Eq a => a -> a -> Bool}
  \item \mintinline{haskell}{h :: a -> a -> Bool}
  \item \mintinline{haskell}{h :: Ord a => a -> a -> String}
  \item \mintinline{haskell}{h :: Show a => a -> a -> String}
  \item \mintinline{haskell}{h :: (Ord a, Show a) => a -> a -> Bool}

\end{enumerate}

\questionTitle{62}{Aninhamento com três casos}
\begin{flushright}
  \footnotesize
  \textit{\textbf{5: Expressão condicional}}\\
  \textit{c05-if-aninhamento-02}\\

\end{flushright}
Avalie a função abaixo na entrada \haskellinline|f 5|:
\begin{haskellcode}
f x = if x > 0
      then if x > 10
           then 1
           else 2
      else 3
\end{haskellcode}

\begin{enumerate}
  \item Ocorre um erro de sintaxe por falta de parênteses.

  \item Ocorre um erro de ambiguidade entre os dois \haskellinline|else|.

  \item \haskellinline|3|

  \item \haskellinline|2|

  \item \haskellinline|1|


\end{enumerate}

\questionTitle{63}{If como subexpressão}
\begin{flushright}
  \footnotesize
  \textit{\textbf{5: Expressão condicional}}\\
  \textit{c05-if-avaliacao-02}\\

\end{flushright}
Considere a expressão:
\begin{haskellcode}
4 * (if 'B' < 'A' then 2 + 3 else 2) - 3
\end{haskellcode}
Qual é o seu valor?

\begin{enumerate}
  \item \haskellinline|17|

  \item \haskellinline|5|

  \item Ocorre um erro de tipos.

  \item \haskellinline|20|

  \item \haskellinline|-1|


\end{enumerate}

\questionTitle{64}{Identificação de violação de pureza (I/O temporal)}
\begin{flushright}
  \footnotesize
  \textit{\textbf{1: Paradigmas de programação}}\\
  \textit{c01-transparencia-referencial-02}\\

\end{flushright}
Considere uma função \haskellinline|obterHoraAtual| que consulta o relógio interno do sistema operacional e retorna uma representação do momento exato em que foi chamada.

Por que a existência desta função viola a propriedade de pureza em um modelo matemático estrito?

\begin{enumerate}
  \item Porque a leitura do relógio do sistema exige chamadas de baixo nível em \emph{Assembly}, que quebram as abstrações do compilador Haskell.

  \item Porque o tempo é uma grandeza contínua e linguagens funcionais só conseguem modelar estruturas de dados e matrizes estritamente discretas.

  \item Ela não viola a pureza; contanto que a função não modifique a hora do sistema (apenas efetue uma leitura), ela continua sendo referencialmente transparente.

  \item Porque ela consome excessivos ciclos de processamento do \emph{garbage collector}, atrasando a execução do fluxo principal da aplicação.

  \item Porque seu valor de retorno depende de um estado global em mutação constante fora do controle do programa, fazendo com que chamadas com os mesmos argumentos produzam resultados diferentes.


\end{enumerate}

\questionTitle{65}{Escopo limitado à equação}
\begin{flushright}
  \footnotesize
  \textit{\textbf{5: Expressão condicional}}\\
  \textit{c05-where-escopo-02}\\

\end{flushright}
Considere a definição:
\begin{haskellcode}
f x
  | y > 0     = y
  | otherwise = 0
  where y = x - 1

f x = y + 10
\end{haskellcode}
Qual afirmação está correta?

\begin{enumerate}
  \item A segunda equação pode usar \haskellinline|y| normalmente, pois \haskellinline|where| tem escopo em toda a função.

  \item A segunda equação pode usar \haskellinline|y|, mas somente se ela também tiver guardas.

  \item Ambas as equações compartilham \haskellinline|y|, mas apenas se houver uma única assinatura de tipo.

  \item A segunda equação não pode usar \haskellinline|y|, pois esse nome vale apenas na primeira equação.

  \item \haskellinline|where| sempre define variáveis globais no módulo.


\end{enumerate}

\questionTitle{66}{Entrada que provoca falha}
\begin{flushright}
  \footnotesize
  \textit{\textbf{5: Expressão condicional}}\\
  \textit{c05-guardas-exaustividade-02}\\

\end{flushright}
Para qual chamada a função abaixo falha em tempo de execução?
\begin{haskellcode}
comparar x y
  | x > y = "Maior"
  | x < y = "Menor"
\end{haskellcode}

\begin{enumerate}
  \item Nenhuma, pois o compilador adiciona \haskellinline|otherwise| implicitamente.

  \item \haskellinline|comparar 5 10|

  \item \haskellinline|comparar 10 5|

  \item \haskellinline|comparar 0 (-1)|

  \item \haskellinline|comparar 5 5|


\end{enumerate}

\questionTitle{67}{Papel de return em ações de E/S recursivas}
\begin{flushright}
  \footnotesize
  \textit{\textbf{10: Ações de E/S recursivas}}\\
  \textit{c10-return-tipos-01}\\

\end{flushright}
Em Haskell, considere a especialização da função \haskellinline|return| usada neste capítulo:

\begin{haskellcode}
return :: a -> IO a
\end{haskellcode}

Qual alternativa descreve corretamente o papel de \haskellinline|return| em uma ação de E/S recursiva?

\begin{enumerate}
  \item Ela imprime automaticamente o valor produzido pela ação de E/S.

  \item Ela transforma um valor puro em uma ação \haskellinline|IO| que, ao ser executada, produz esse valor sem realizar interação externa.

  \item Ela lê um novo valor da entrada padrão e o devolve ao próximo passo da recursão.

  \item Ela interrompe imediatamente a execução da ação atual, como o comando \haskellinline|return| em C, Java ou Python.

  \item Ela reinicia a ação recursiva, funcionando como um comando \haskellinline|continue|.


\end{enumerate}

\questionTitle{68}{A função putStr}
\begin{flushright}
  \footnotesize
  \textit{\textbf{8: Programas interativos}}\\
  \textit{c08-entrada-saida-01}\\

\end{flushright}
Qual é o tipo de \haskellinline|putStr| e qual é seu comportamento?

\begin{enumerate}
  \item \haskellinline|String -> IO String|. Ela imprime a string e também a retorna para o programa.

  \item \haskellinline|String -> IO ()|. Ela constrói uma ação que imprime a string na saída padrão sem acrescentar uma quebra de linha.

  \item \haskellinline|String -> String|. Ela apenas converte uma string em outra string.

  \item \haskellinline|IO String|. Ela lê uma string da entrada padrão.

  \item \haskellinline|IO ()|. Ela imprime uma linha vazia sem receber argumentos.


\end{enumerate}

\questionTitle{69}{Divisão inteira por subtrações sucessivas}
\begin{flushright}
  \footnotesize
  \textit{\textbf{9: Funções recursivas}}\\
  \textit{c09-projeto-divisao-inteira-01}\\

\end{flushright}
Para \haskellinline|p >= 0| e \haskellinline|q > 0|, qual ideia recursiva calcula corretamente o quociente da divisão inteira de \haskellinline|p| por \haskellinline|q|?

\begin{enumerate}
  \item Usar o caso base \haskellinline|q == 0| para retornar o quociente.

  \item Enquanto \haskellinline|p >= q|, subtrair \haskellinline|q| de \haskellinline|p| e somar \haskellinline|1| ao quociente.

  \item Enquanto \haskellinline|p >= q|, somar \haskellinline|q| a \haskellinline|p| e subtrair \haskellinline|1| do quociente.

  \item Retornar sempre \haskellinline|p| quando \haskellinline|p < q|.

  \item Dividir \haskellinline|q| por \haskellinline|p| em cada chamada recursiva.


\end{enumerate}


\end{multicols}
\makeanswersheet{UFOP}{BCC222: Programação Funcional}{José Romildo}{2026/1}{Primeira Prova}{10}{69}{25}{7}
\gabarito{10}{B,B,D,E,B,A,D,A,B,A,C,C,E,C,A,E,A,D,E,D,C,A,A,E,E,B,B,E,A,A,E,D,D,A,D,A,B,E,D,E,C,C,B,E,C,A,E,B,A,D,D,B,C,E,B,C,B,A,C,A,E,D,B,E,D,E,B,B,B}
\end{document}
